% In C. Verrem Actio Secunda IV --- headnote
% in-verrem-2-4, autumn 70 BC

\textit{Book IV of the five-book Actio Secunda --- ``de signis,'' on
Verres's plundering of the art and sacred objects of Sicily. The
book is the most famous of the Verrines for its set-piece
narrative passages: the Heius collection at Messana (the four
statues from the household shrine, including the Praxitelean
Cupid); the Cibyratan brothers Hiero and Tlepolemus who tracked
out plate ``like hunting-dogs'' across the province; the
candelabrum of Antiochus, set aside for the Capitoline and
diverted into Verres's house; Diodorus the Maltese threatened
with a capital charge over a set of Mentor cups; the despoiling
of the Diana of Segesta (Scipio's restitution after the third
Punic War, the women of the city anointing her with unguents and
crowns as she was carried out to the sea); the Mercury of
Tyndaris and the binding of the proagorus Sopater, naked, in
winter rain, to the equestrian statue of Marcellus; the temple
of Hercules at Agrigentum, whose plundering by night raised
the whole city in arms; the temple of Ceres at Henna, where
Verres became ``a second Orcus'' carrying off Ceres herself;
and the long, set-piece comparison of the two takings of
Syracuse --- by Marcellus, who spared the temples; and by
Verres, who stripped them. Out of the temple of Minerva on
Ortygia he took the cavalry-battle panels, the gold studs from
the doors, the Gorgon's head, the Sappho of Silanion from the
prytaneum, the great image of Imperator Jupiter (the third of
that triad of statues, the others being Flamininus's on the
Capitoline and the one at the entrance to the Pontus). The
guides who used to show what was where now ``show what has been
taken away from where.'' The book closes on the Syracusan senate
revoking by acclamation the praise of Verres they had been
forced to issue, and the Mamertine praise (the only one
remaining in the province) collapsing under questioning.}
